%% This file contains a number of tweaks that are typically applied to the main document.
%% They are not enabled by default, but can be enabled by uncommenting the relevant lines.

%%
%% Inline annotations; for predefined colors, refer to "dvipsnames" in the xcolor package:
%% https://tinyurl.com/overleaf-colors
%%
\newcommand{\red}[1]{{\color{red}#1}}
\newcommand{\todo}[1]{{\color{red}#1}}
\newcommand{\TODO}[1]{{\color{gray}#1}}  % gray hint: what to write here
%%
%% disable for camera ready / submission by uncommenting these lines  
%%
% \renewcommand{\TODO}[1]{}
% \renewcommand{\todo}[1]{#1}

%%
%% work harder in optimizing text layout. Typically shrinks text by 1/6 of page, enable
%% it at the very end of the writing process, when you are just above the page limit
%%
% \usepackage{microtype}

%%
%% fine-tune paragraph spacing
%%
% \renewcommand{\paragraph}[1]{\vspace{.5em}\noindent\textbf{#1.}}

%%
%% globally adjusts space between figure and caption
%%
% \setlength{\abovecaptionskip}{.5em}


%%
%% Allows "the use of \paper to refer to the project name"
%% with automatic management of space at the end of the word
%%
% \usepackage{xspace}
% \newcommand{\paper}{ProjectName\xspace}

%%
%% Commonly used math definitions
%%
% \DeclareMathOperator*{\argmin}{arg\,min}
% \DeclareMathOperator*{\argmax}{arg\,max}

%%
%% Tigthen underline
%%
% \usepackage{soul}
% \setuldepth{foobar}

%%
%% Lab3 additions
%%
\usepackage{array}   % m{} columns: vertically centered image grids
\usepackage{listings}
\lstset{language=Python, basicstyle=\ttfamily\scriptsize, columns=fullflexible,
        breaklines=true, frame=single, keywordstyle=\color{NavyBlue}\bfseries,
        commentstyle=\color{gray}, showstringspaces=false}

%% \img[width]{path}: a framed placeholder box (path label scaled to fit) until the image file exists;
%% raised to sit on the baseline like a real image, so row labels center the same way
\newcommand{\img}[2][\linewidth]{\IfFileExists{#2}{\includegraphics[width=#1]{#2}}{%
  \raisebox{\depth}{\fbox{\parbox[c][\dimexpr#1-2\fboxsep-2\fboxrule][c]{\dimexpr#1-2\fboxsep-2\fboxrule}{\centering
  \resizebox{\linewidth}{!}{\parbox{2cm}{\centering\tiny\def\UrlBreaks{\do\/\do\_}\nolinkurl{#2}}}}}}}}
\newcolumntype{C}[1]{>{\centering\arraybackslash}m{#1}}

%% Images live in Figure/: Figure/{sds,vsd,sdi}/<prompt>.png, Figure/ablation/<setting>/<value>/<prompt>.png
%% \resrow{method dir}{row label}: one row with all 10 evaluation prompts
\newcommand{\resrow}[2]{\rotatebox{90}{\scriptsize #2}
  & \img{Figure/#1/A_red_bus_driving_on_a_desert_road.png} & \img{Figure/#1/a_boat_in_a_river.png}
  & \img{Figure/#1/A_cabin_surrounded_by_forests.png}      & \img{Figure/#1/A_church_beside_a_lake.png}
  & \img{Figure/#1/A_villa_close_to_the_pool.png}          & \img{Figure/#1/A_castle_next_to_a_river.png}
  & \img{Figure/#1/A_burger_on_the_table.png}              & \img{Figure/#1/A_dog_sitting_on_grass.png}
  & \img{Figure/#1/a_cat_sitting_on_a_table.png}           & \img{Figure/#1/A_car_on_the_road.png} \\}

%% \ablpanel{setting}{5 values}: value labels on top, then 3 labeled prompts (rows) x 5 values (columns)
%% \ablcell: ablation image size; Fig. 6 uses 1.5\ablcell, so its 2 rows are as tall as the 3 ablation rows
\newlength{\ablcell}\setlength{\ablcell}{.087\textwidth}
\newcommand{\ablrow}[8]{\rotatebox{90}{\scriptsize #3} & \img{Figure/ablation/#1/#4/#2.png} & \img{Figure/ablation/#1/#5/#2.png}
  & \img{Figure/ablation/#1/#6/#2.png} & \img{Figure/ablation/#1/#7/#2.png} & \img{Figure/ablation/#1/#8/#2.png} \\}
\newcommand{\ablpanel}[6]{%
  \begin{tabular}{@{}m{1em}*{5}{C{\ablcell}}@{}}
    & \scriptsize #2 & \scriptsize #3 & \scriptsize #4 & \scriptsize #5 & \scriptsize #6 \\
    \ablrow{#1}{A_burger_on_the_table}{Burger}{#2}{#3}{#4}{#5}{#6}
    \ablrow{#1}{a_cat_sitting_on_a_table}{Cat}{#2}{#3}{#4}{#5}{#6}
    \ablrow{#1}{A_villa_close_to_the_pool}{Villa}{#2}{#3}{#4}{#5}{#6}
  \end{tabular}}
